% Exact worked example: a2=(1,2); q1=(3/5,4/5); projection=(33/25,44/25).
% The remainder (-8/25,6/25) is perpendicular to q1, length2/5.
\begin{tikzpicture}[>=Latex,font=\small,x=1.4cm,y=1.4cm]
\draw[->,black!40] (-.3,0)--(2.1,0);
\draw[->,black!40] (0,-.15)--(0,2.3);
\draw[->,very thick,blue!65!black] (0,0)--(1,2) node[above left] {$\mathbf a_2$};
\draw[->,thick,dashed] (0,0)--(1.32,1.76);
\node[anchor=west] at (1.39,1.65) {$r_{12}\mathbf q_1$};
\draw[->,very thick,orange!80!black] (1.32,1.76)--(1,2);
\draw[->,thick] (0,0)--(.6,.8) node[right] {$\mathbf q_1$};
\draw[thin] (1.32-.12,1.76-.16)--(1.32-.28,1.76-.04)--(1.32-.16,1.76+.12);
\node[align=left,anchor=west] at (3.5,1.7) {$r_{12}=\mathbf q_1^{\mathsf T}\mathbf a_2$\\[5pt]
 $\mathbf v=\mathbf a_2-r_{12}\mathbf q_1$\\[5pt]
 $\mathbf q_1^{\mathsf T}\mathbf v=0$};
\node[align=left,anchor=west] at (3.5,.3) {normalize the perpendicular\\remainder to obtain $\mathbf q_2$};
\node[anchor=west] at (1.6,2.4) {$\mathbf v$};
\draw[thin] (1.6,2.3)--(1.12,1.94);
\end{tikzpicture}
